\chapter{Reflow and resize}
\label{ch:reflow}

\begin{aside}
The terminal resizes. Your UI must follow. Layout knows how;
the trick is doing it fast enough that drag-resize feels live.
\end{aside}

\section{The signal}

SIGWINCH fires when the terminal size changes. \maya{}'s
loop catches the EINTR return, queries the new size via
\code{TIOCGWINSZ}, and triggers a resize event.

\section{The reflow}

A resize event invalidates both buffers:

\begin{enumerate}[leftmargin=*]
  \item Reallocate canvases to the new dimensions.
  \item Fill the front with sentinels.
  \item Re-run layout with the new viewport.
  \item Repaint and full-screen emit.
\end{enumerate}

The cost: one allocation pair (front + back), one layout
pass, one full-screen ANSI emission. On a 200x60 terminal,
all under 5 ms.

\section{Layout responsiveness}

If layout is slow (deep nesting, expensive measure
functions), drag-resize lags. The user drags the corner and
watches stale frames before the new size catches up.

\maya{}'s layout is fast enough for typical UIs that
drag-resize is live.

\section{Coalescing rapid resizes}

A user dragging the resize handle generates many SIGWINCH
events. Each is a few bytes of ioctl work but a full
relayout is expensive. \maya{} coalesces: if multiple
resize events arrive in one wakeup, only the latest size
triggers a relayout.

\section{Reflow vs scroll}

For a scrolling region, resize changes the viewport. The
content's logical scroll position is preserved; the visible
rows shift.

For a wrapping text view, the wrap points change. Cursor
position in characters is preserved; cursor position in
visual rows shifts. The view scrolls so the cursor stays
visible.

\section{Sticky pinned content}

Status bars, headers, footers should not move. Layout
naturally handles this: they have fixed heights and live in
their own row of a column layout.

\section{Content-driven sizing}

A widget that auto-sizes to its content (a label) has zero
flex --- its size is dictated by content. Reflow recomputes
its measure; if content fits, no change. If content
overflows, fall back to wrap or truncate.

\section{The minimum viable size}

A 1x1 terminal is technically possible. Reasonable UIs
expect at least 20x5; very rich UIs require 80x24 or more.

\maya{}'s elements have minimum sizes (border = 3x3, button
= label-width + 2, etc.). Below these, layout warns and
falls back to a degraded form (no border, ASCII fallback).

\section{Resize event payload}

The event carries the new size:

\begin{lstlisting}
struct ResizeEvent {
    Size old_size;
    Size new_size;
};
\end{lstlisting}

Widgets can subscribe and react (e.g., recompute their
internal scroll position).

\section{Recap}

\begin{recap}
\begin{itemize}[leftmargin=*]
  \item SIGWINCH + TIOCGWINSZ + sentinel flood = resize.
  \item Coalesce burst events.
  \item Layout speed = drag-resize feel.
  \item Status bars stay put via layout, not magic.
  \item Mind the minimum viable size.
\end{itemize}
\end{recap}

\section*{Workshop}

\begin{enumerate}[leftmargin=*]
  \item Drag-resize your terminal during a \maya{} app.
    Watch the relayout happen smoothly.
  \item Print layout time on resize. Identify any deep
    tree that pushes you over 5 ms.
  \item Build a widget that displays current size in its
    own corner.
\end{enumerate}
